On the Einstein region with $\Lambda \neq 0$ the frozen point $I$ is the parallel scale tractor, of squared norm $\ip{I}{I} = \Lambda/3 = H^{2}$, and in the flat twistor chart its frozen point is $I_{\Lambda} = n_{\infty} + (\Lambda/6)\,n_{0}$. Mass is conformal weight read against it: a density of weight $w$ coupled to the frozen point obeys the Klein-Gordon equation with
\[
m^{2} = -\,w(w+3)\,\frac{\Lambda}{3} = \Delta(3-\Delta)\,H^{2}, \qquad \Delta = -w,
\]
so a real weight gives $m^{2} \ge 0$ only on $w \in [-3, 0]$, where $m \le 3H/2$, and every heavier field has complex weight $w = -\tfrac32 \pm i\nu$ with $\nu \approx m/H$. In particular the weight of the Planck mass is complex, and its imaginary part is
\[
N := \operatorname{Im} w(M_{P}) \approx M_{P}/H, \qquad\text{of order } 10^{60}.
\]
Throughout this paper $N$ means $M_{P}/H$, the number the dictionary attaches to the Planck mass. \emph{Provenance:} \cite[Thms.~4.5, 4.7]{ButlerT6} for the frozen point and its norm; \cite[Thm.~3.1, Cor.~3.2]{ButlerT11} for the dictionary and the complex weights; \cite[Lem.~5.1]{ButlerT11} for the chart form of $I_{\Lambda}$. The value $10^{60}$ is the measured ratio $M_{P}/H$ and carries the record's status \cite{PDG2024}.
